The process of \betareduction is somehow like calculating in arithmetic:
\[
    (3 + 4) \times (4 - 2) + 1 \quad\to\quad 7 \times (4 - 2) + 1 \quad\to\quad 7 \times 2 + 1 \quad\to\quad 14 + 1 \quad\to\quad 15
\]
Note that one could have also done it in a different order, that is:
\[
    (3 + 4) \times (4 - 2) + 1 \quad\to\quad 7 \times 2 + 1 \quad\to\quad 7 \times (4 - 2) + 1 \quad\to\quad 14 + 1 \quad\to\quad 15
\]
However, either way, both calculations give the same result, the normal form of $15$. Obviously every expression above is equivalent to $15$ but we usually refer to $15$ because its syntax is less verbose.

With \taureduction, we also introduce some kind of calculation where we eliminate the bound variables in the body of an expression until there is no more to eliminate.

\begin{reusable}{definition}{definitionTauRedex}
    An expression $M$ is a \tauredex if either:
    \begin{itemize}
        \item $M$ is a \betaredex.
        \item $\text{BV}(M) \ne \emptyset$
    \end{itemize}
\end{reusable}

\begin{reusable}{definition}{definitionTauNormalForm}
    We say that $M$ is a \taunormalform, equivalently written $\taunf$, if $M$ does not contain any \tauredex, that is, if $M$ cannot be reduced any further. We also say that $M$ is in $\taunf$.
\end{reusable}

\begin{reusable}{definition}{definitionTauHasNormalForm}
    We say that $M$ has a $\taunf$ if there is an $N$ in \taunormalform such that $M \taueq N$.
\end{reusable}

Obviously, if a $M$ is in $\taunf$ and reduces to $N$ then $N$ is necessarily \alphaequivalent to $M$.

\begin{reusableWithProof}{lemma}{lemmaTauReductionOfTauNfImpliesAlphaEqWithContractum}
    If $M$ is in $\taunf$ and $M \totau N$, then we have $M \alphaeq N$.
\end{reusableWithProof}

\begin{reusable}{lemma}{lemmaThereExistTermsWithNoNormalForm}
    In $\LambdaTau$, there are terms that do not have a \taunormalform.
\end{reusable}

\begin{proof}
     $\LambdaTau$ is an extension of pure \lambdacalc, that is, $\Lambda \subseteq \LambdaTau$ and therefore all expressions in pure \lambdacalc are valid in $\LambdaTau$. The fixed point theorem proven by Haskell and Curry made them build the Y combinator term that is defined as follows: For any $L$, $Y:=(\lambda x.L(xx))(\lambda x.L(xx))$. This expression \betareduces to itself. Since \betareduction implies \taureduction, it \taureduces to itself and still contains a \tauredex. Hence, it never ends up in a state where it has no \tauredex. We conclude that it has no \taunormalform.
\end{proof}

\begin{reusable}{definition}{definitionFiniteReductionPath}
    A \emph{finite reduction path} from \(M\) is a finite sequence of terms $N_0, N_1, N_2, \ldots, N_n$ such that \(N_0 \alphaeq M\) and \(N_i \totau N_{i+1}\) for each \(i\) with \(0 \le i < n\).
\end{reusable}

\begin{reusable}{definition}{definitionInfiniteReductionPath}
    An \emph{infinite reduction path} from \(M\) is an infinite sequence $N_0, N_1, N_2, \ldots$ with \(N_0 \alphaeq M\) and \(N_i \totau N_{i+1}\) for all \(i \in \mathbb{N}\).
\end{reusable}

\begin{reusable}{definition}{definitionNormalisation}[ weak and strong normalisation ]
    \begin{enumerate}
        \item \( M \) is \textit{weakly normalising} if there exists an \( N \) in \taunormalform such that \( M \totau N \).
        \item \( M \) is \textit{strongly normalising} if there are no infinite reduction paths starting from \( M \).
    \end{enumerate}
\end{reusable}

At this stage we want to prove that $\LambdaTau$ has the diamond property, that is, if $M \totau N_{1}$ and $M \totau N_{2}$, then there exists $N_{1} \totau N$ and $N_{2} \totau N$. However, to prove that proposition more easily we shall introduce the concept of parallel \taureduction to avoid messy commutation proofs.

\begin{reusable}{definition}{definitionParallelTauReduction}[ parallel \taureduction]
    \begin{align}
        a \in \mathbb{V} & \vdash a \toptau a \\
        A \toptau A' \land N \toptau N' & \vdash (\lambda x.A)N \toptau A'[x := N'] \\
        A \toptau A' \land N \toptau N' \land M \toptau M' \land M' \alphaeq N' & \vdash (\lambda x \oftype M.A)N \toptau A'[x := N'] \\
        M \toptau M' \land N \toptau N' & \vdash \lambda x \oftype M.N \toptau \lambda x \oftype M'.N'[x := M'] \\
        M \toptau M' & \vdash  \lambda x.M \toptau \lambda x.M' \\
        M \toptau M' \land N \toptau N' & \vdash \lambda x \oftype M.N \toptau \lambda x \oftype M'.N' \\
        M \toptau M' \land N \toptau N' & \vdash MN \toptau M'N'
    \end{align}
\end{reusable}

\begin{reusable}{theorem}{theoremDiamondPropertyParallel}
    If $M \toptau N_{1}\text{ and }M \toptau N_{2}$, then there exists a $N$ such that: $N_{1} \toptau N \text{ and }N_{2} \toptau N$.
\end{reusable}

\begin{proof}
We prove the property by structural induction on the parallel \taureduction rules.

\paragraph{Base Case}
Suppose $M$ is a variable $x$. By applying the base rule we have $x \toptau x$ and necessarily $N_1=N_2=x$. If we take $N=x$, we have $N_1 \toptau N$ and $N_2 \toptau N$.

\paragraph{Inductive Case: Untyped abstraction compatibility}
Suppose $M=\lambda x.N$. If $M \toptau M_1$ and $M \toptau M_2$, then there exists $N_1$ and $N_2$ such that $M_1=\lambda x.N_1$ and $M_2=\lambda x.N_2$. However, by the induction hypothesis, we know that there is a $P'$ such that $N_1 \toptau P'$ and $N_2 \toptau P'$. We apply the untyped abstraction compatibility rule and we get $M_1=\lambda x.N_1 \toptau \lambda x.P'$ and $M_2=\lambda x.N_2 \toptau \lambda x.P'$. We conclude that $\lambda x.P'$ is our $N$ and therefore $M_1 \toptau N$ and $M_2 \toptau N$.

\paragraph{Inductive Case: Typed abstraction compatibility}
There are two cases: either the variable of the abstraction is in the bound variables of the body or it is not. Let's start with the case where it is not.

\subparagraph{Variable not in Bound Variables}
Suppose $M=\lambda x \oftype T.N$. If $M \toptau M_1$ and $M \toptau M_2$, then there exists $N_1, T_1, N_2, T_2$ such that $M_1:=\lambda x \oftype T_1.N_1$ and $M_2:=\lambda x \oftype T_2.N_2$ with $T \toptau T_1$ and $T \toptau T_2$. However, by the induction hypothesis, we know that there is a $P'$ such that $N_1 \toptau P'$ and $N_2 \toptau P'$ and a $T'$ such that $T_1 \toptau T'$ and $T_2 \toptau T'$. We apply the typed abstraction body (rule 6). rule and we get $M_1=\lambda x \oftype T_1.N_1 \toptau \lambda x \oftype T'.P'$ and $M_2=\lambda x \oftype T_2.N_2 \toptau \lambda x \oftype T'.P'$. We conclude that $\lambda x \oftype T'.P'$ is our $N$ and therefore $M_1 \toptau N$ and $M_2 \toptau N$.

\subparagraph{Variable in Bound Variables}
    Suppose $M=\lambda x \oftype T.N[x := T]$. If $M \toptau M_1$ and $M \toptau M_2$, then there exists $T_1', N_1', T_2', N_2'$ such that:

    \[
    \begin{array}{l}
        M_1:=\lambda x \oftype T_1'.N_1'[x := T_1'] \\
        M_2:=\lambda x \oftype T_2'.N_2'[x := T_2'] \\
        T \toptau T_1' \text{ and } T \toptau T_2' \\
        x \in \BV{N_1'} \text{ and } x \in \BV{N_2'}
    \end{array}
    \]
    
    By the induction hypothesis we know that there exist $T'$ and $N'$ such that:
    \[
    \begin{array}{l}
        T_1' \toptau T' \\
        T_2' \toptau T' \\
        N_1' \toptau N' \\
        N_2' \toptau N'
    \end{array}
    \]
    
    We apply the rule 4 on $T_1', N_1'$ because $T_1' \toptau T'$, $N_2' \toptau N'$ and $x \in \BV{N_1'}$. We deduce $\lambda x \oftype T_1'.N_1' \toptau \lambda x \oftype T_1'.N_1'[x := T_1']$. The same applies with $T_2'$ and $N_2'$ and we therefore also deduce $\lambda x \oftype T_2'.N_2' \toptau \lambda x \oftype T_2'.N_2'[x := T_2']$. Those two results means that

\paragraph{Inductive Case: Application compatibility}
    
\end{proof}

\begin{reusable}{theorem}{theoremConfluence}[Confluence theorem]\\
Let $M, N_1, N_2$ such that $\quad M \totau N_{1} \quad\text{and}\quad M \totau N_{2}$.\\
Then there is a $N$ such that: $N_{1} \totau N \quad\text{and}\quad N_{2} \totau N$.
\end{reusable}
